\documentclass[13pt,letter]{article}

\usepackage[margin=1in]{geometry}
\usepackage{hyperref}
\usepackage{graphicx}
\usepackage{longtable}
\usepackage{booktabs}
\usepackage{listings}
\usepackage{enumitem}
\usepackage{setspace}
\usepackage{textcomp} % fixes \text related errors
\usepackage{amsmath}

\title{Standard Operating Procedure (SOP)\\
Whole Genome NGS Pipeline (wf\_ngs\_pipeline.wdl)}
\author{Prepared by: Dieter Best}
\date{\today}

\begin{document}

\maketitle
\tableofcontents
\newpage

\section{Purpose}

This SOP describes how to run and interpret results from the \texttt{wf\_ngs\_pipeline.wdl} workflow from the GitHub repository:

\begin{itemize}
  \item \url{https://github.com/ranlib/whole_genome}
\end{itemize}

The workflow processes paired-end Illumina whole-genome sequencing (WGS) data and performs:

\begin{itemize}
  \item Quality control and read trimming
  \item Taxonomic classification
  \item Reference alignment
  \item Coverage and QC metrics
  \item Variant calling and annotation
  \item Summary reporting with MultiQC
\end{itemize}

\section{Pipeline Overview}

The pipeline is written in Workflow Description Language (WDL) and processes multiple samples in parallel. Each sample passes through the following stages:

\begin{enumerate}
  \item Raw read QC
  \item Read trimming and filtering
  \item Post-trimming QC
  \item Taxonomic classification
  \item Alignment to reference genome
  \item BAM metrics and whole-genome QC
  \item Coverage estimation
  \item Variant calling and structural variant detection
  \item Variant annotation
  \item Aggregated reporting
\end{enumerate}

\section{Workflow Steps}

\subsection{Pre-Alignment Quality Control}

\begin{itemize}
  \item \textbf{SeqKit}: Summary statistics of raw reads.
  \item \textbf{FastQC}: Quality metrics on raw FASTQ files.
  \item \textbf{Fastp (Loose \& Tight modes)}: Adapter trimming and quality filtering.
  \item \textbf{SeqKit (post-cleaning)}: Statistics after trimming.
  \item \textbf{FastQC (post-cleaning)}: QC reports on cleaned reads.
\end{itemize}

\subsection{Taxonomic Classification}

\begin{itemize}
  \item \textbf{Centrifuge}: Classification of reads to identify contamination or microbial composition.
\end{itemize}

\subsection{Alignment and BAM Processing}

\begin{itemize}
  \item \textbf{Minimap2}: Alignment of reads to reference genomes.
  \item \textbf{SAMtools}: Indexing reference FASTA and BAM files.
  \item \textbf{Picard / SAMtools metrics}: Alignment and insert-size metrics.
\end{itemize}

\subsection{Coverage and Whole-Genome QC}

\begin{itemize}
  \item \textbf{GATK WGS Metrics}: Coverage and quality metrics.
  \item \textbf{Mosdepth}: Depth of coverage statistics.
\end{itemize}

\subsection{Variant Calling and Annotation}

\begin{itemize}
  \item \textbf{GATK}: Variant calling producing VCF files.
  \item \textbf{Delly}: Structural variant detection.
  \item \textbf{VCF concatenation}: Merge SNVs and SVs.
  \item \textbf{SnpEff}: Variant effect annotation.
\end{itemize}

\subsection{Summary Reporting}

\begin{itemize}
  \item \textbf{MultiQC}: Aggregates all QC, alignment, coverage, and variant reports into a single HTML report.
\end{itemize}

\section{Inputs}

\begin{longtable}{ll}
\toprule
\textbf{Parameter} & \textbf{Description} \\
\midrule
reads1 & Forward FASTQ files \\
reads2 & Reverse FASTQ files \\
samplenames & Unique sample identifiers \\
references & Reference genome FASTA files \\
dockerImages & Map of Docker images for tools \\
threads & Number of CPU threads \\
memory & Memory allocation per task \\
disk\_size & Disk allocation for large tasks \\
\bottomrule
\end{longtable}

\section{Outputs}

\begin{longtable}{ll}
\toprule
\textbf{Output} & \textbf{File Type} \\
\midrule
Cleaned reads & FASTQ \\
FastQC reports & HTML, ZIP \\
Centrifuge classification & TSV \\
Alignments & BAM, BAI \\
Coverage reports & TXT, TSV \\
Variant calls & VCF, VCF.GZ \\
Annotated variants & VCF \\
MultiQC report & HTML \\
\bottomrule
\end{longtable}

\section{Execution Instructions}

Example command using Cromwell:

\begin{lstlisting}[language=bash]
  java -jar cromwell.jar run wf_ngs_pipeline.wdl -i inputs.json
\end{lstlisting}

\section{Best Practices}

\begin{itemize}
  \item Validate FASTQ file integrity before running the pipeline.
  \item Ensure reference genomes are properly indexed.
  \item Verify Docker images are accessible.
  \item Adjust memory and CPU resources based on dataset size.
  \item Review MultiQC output for global QC assessment.
\end{itemize}

\section{Troubleshooting}

\begin{itemize}
  \item \textbf{Workflow fails at startup}: Check all required inputs are provided in \texttt{inputs.json}.
  \item \textbf{Out-of-memory errors}: Increase memory settings for alignment or variant calling tasks.
  \item \textbf{Empty MultiQC report}: Ensure all upstream tasks completed successfully.
\end{itemize}

\section{References}

\begin{itemize}
  \item \url{https://github.com/ranlib/whole_genome}
  \item \url{https://openwdl.org}
  \item \url{https://github.com/broadinstitute/cromwell}
\end{itemize}

\end{document}
